% XiaoHongShu-style long-note example
\documentclass[12pt]{article}
\usepackage{xhsnote-template}

% --------------------------------
% Optional custom macros
% --------------------------------
\newcommand{\E}{\mathbb{E}}
\newcommand{\R}{\mathbb{R}}
\newcommand{\dd}{\,\mathrm{d}}
\newcommand{\norm}[1]{\left\lVert #1 \right\rVert}

\begin{document}

\MakeShortHeader
  {I-MMSE Identity}
  {A compact note on mutual information and MMSE in Gaussian channels}
  {Information Theory, Gaussian Channel, MMSE}

\begin{KeyBox}
在高斯信道
\[
Y_\gamma=\sqrt{\gamma}X+Z,\qquad Z\sim N(0,I_d),\quad Z\perp X
\]
中，互信息关于信噪比 \(\gamma\) 的导数恰好等于 MMSE：
\[
\frac{\dd}{\dd\gamma}I(X;Y_\gamma)=\frac{\log e}{2}\,\mathrm{mmse}(X\mid Y_\gamma).
\]
\end{KeyBox}

\section{问题背景}

I-MMSE identity 是信息论与统计估计之间一个非常漂亮的桥梁。
它说明：\Term{信息增益的变化率}，正好由\Term{最优估计误差}刻画。

\begin{XHSDefinition}[MMSE]
设 \(X\in\R^d\) 且 \(\E[\norm{X}^2]<\infty\)。定义
\[
\mathrm{mmse}(X\mid Y):=\E\!\left[\norm{X-\E[X\mid Y]}^2\right].
\]
\end{XHSDefinition}

\begin{XHSTheorem}[I-MMSE]
设 \(Y_\gamma=\sqrt{\gamma}X+Z\)，其中 \(Z\sim N(0,I_d)\) 与 \(X\) 独立，则
\[
\frac{\dd}{\dd\gamma}I(X;Y_\gamma)=\frac{\log e}{2}\,\mathrm{mmse}(X\mid Y_\gamma).
\]
\end{XHSTheorem}

\begin{xhsFormulaBox}[一个常见推论]
由 I-MMSE 可得
\[
I(X;Y_\gamma)=\frac{\log e}{2}\int_0^\gamma \mathrm{mmse}(X\mid Y_t)\,\dd t.
\]
\end{xhsFormulaBox}

\begin{Takeaway}
在 Gaussian noise 下，“知道多少”和“还能估错多少”本质上是同一个结构的两种表达。
\end{Takeaway}

\end{document}
